%============================================================================================================
%DISTANCE DE DECOLLAGE AVEC LES DIFFERENTES METHODE 
%================================================================================================================
%------------------------------------------------
\section{\textcolor{blue}{Distance parcourue au sol} }

En appliquant la seconde loi de Newton à l’avion pendant sa phase de roulage au sol, on obtient :

\begin{equation}
	\frac{W}{g}\frac{dV}{dt}=T-D-F_r
	\tag{1.1}
\end{equation}

Le principe fondamental de la dynamique (PFD) est utilisé de différentes manières dans le calcul de la distance de décollage.
	
	%------------------------------------------------
	\subsection{\textcolor{green} {Modèle 1 : méthode analytique} }
	Cette méthode est présentée dans \cite{War:09}. 
	\textbf{\textcolor{yellow}{Avantages}}
	
	Cette méthode est applicable à tous les avions et les forces dépendent de la vitesse et sont implementé en focntion de la vitesse a chaque instant. Elle peut être utilisée pour les avions à réaction ainsi que pour les avions à hélices.
	
	\vspace{0.3cm}
	
	\textbf{\textcolor{yellow}{Aproximation}}
	
	.............................
	
	\vspace{0.3cm}
	
	
	En utilisant la relation de dérivation totale :
	
	\begin{equation}
		\frac{dV}{dt}=\frac{ds}{dt}\frac{dV}{ds}=(V-V_{hw})\frac{dV}{ds}
	\end{equation}
	
	On obtient alors :
	
	\begin{equation}
		\frac{ds}{dV}=\frac{W(V-V_{hw})}{g(T-D-F_r)}
	\end{equation}
	
	En intégrant entre deux vitesses successives :
	
	\begin{align}
		S_{i+1}-S_i= \frac{1}{g} \int_{V_i}^{V_{i+1}} \frac{V-V_{hw}} {\dfrac{T-D-F_r}{W}} \, dV
	\end{align}
	
	Dans ce modèle, les forces sont de la forme :
	
	\begin{align}
		D &= \frac{1}{2}\rho V^2 S_w C_D \\
		F_r &= \mu_r W - \frac{1}{2}\mu_r \rho S_w C_L V^2 \\
		T &= (T_0)_i + (T')_i V + (T'')_i V^2
	\end{align}
	
	En posant :
	
	\begin{align}
		\frac{T-D-F_r}{W}=A+BV+CV^2
	\end{align}
	
	avec :
	
	\begin{align}
		A &= \frac{(T_0)_i}{W}-\mu_r \\
		B &= \frac{(T')_i}{W} \\
		C &= \frac{(T'')_i}{W} 	-\frac{1}{2} \frac{\rho}{W/S_w} (\mu_r C_L-C_D)
	\end{align}
	
	L’équation devient :
	
	\begin{align}
		S_{i+1}-S_i	=	\frac{W}{g}	\int_{V_i}^{V_{i+1}} \frac{V-V_{hw}} {A+BV+CV^2} \, dV
	\end{align}
	
	\vspace{1cm}
	
	%------------------------------------------------
	\subsubsection{\textcolor{red}{Modélisation des forces}}
	
	\textbf{Traînée}
	
	La force de traînée est donnée par :
	
	\begin{align}
		D=\frac{1}{2}\rho V^2 S_w C_D
	\end{align}
	
	avec
	
	\begin{align}
		C_D= C_{D0}	+	C_{D0,L}C_L +
		\frac{	\left(\dfrac{16h_w}{b_w}\right)^2  }{ 1+\left(\dfrac{16h_w}{b_w}\right)^2 }	\cdot \frac{C_L^2}{\pi e R_A}
	\end{align}
	
	où :
	
	\[	h_w : \text{ hauteur de l’aile par rapport au sol}	\]	
	\[ b_w :	\text{ envergure de l’aile}	\]
	
	\vspace{0.5cm}
	
	\textbf{Force de frottement}
	
	Pendant le roulage :
	
	\begin{equation}
		F_r= \mu_r(W-L)	= \mu_r	\left( 	W-\frac{1}{2}\rho V^2 S_w C_L  \right)
	\end{equation}
	
	\vspace{0.5cm}
	
	\textbf{Poussée}
	
	La poussée est modélisée par :
	
	\begin{equation}
		T=(T_0)_i+(T')_iV+(T'')_iV^2
	\end{equation}
	
	où $(T_0)_i$, $(T')_i$ et $(T'')_i$ sont des coefficients déterminés expérimentalement.

%------------------------------------------------
	\subsection{\textcolor{blue} { Modèle 2 : méthode simplifiée et méthode numérique   }}
	Cette méthode est présentée dans \cite{smi:60}.
%============================================================================================================
%
%================================================================================================================
		%------------------------------------------------
		\subsubsection{\textcolor{green} { Méthode générale d’estimation utilisant une vitesse moyenne de décollage } }
		
		\textbf{\textcolor{yellow}{Avantages}}
		
		Cette méthode est applicable à tous les avions si les forces peuvent être approximées. Elle peut être utilisée pour les avions à réaction ainsi que pour les avions à hélices.
		
		\vspace{0.3cm}
		
		\textbf{\textcolor{yellow}{Aproximation}}
		
		On considère les forces constantes pendant le roulage sur la piste. Elles sont calculées à la vitesse moyenne $V_R/2$, qui permet une bonne approximation sur cet intervalle.
		
		\vspace{0.3cm}
		
		\textbf{Méthode}
		
		La distance de décollage après application du PFD est donnée par :
		
		\begin{equation*}
			S_G= \frac{ V_R^2 W	}{ 2g[T-D-\mu(W-L)]	}
			\quad
			\text{évaluée en }
			\left(\frac{V_R}{2}\right)
		\end{equation*}
		
		La poussée dépend ici du type de moteur.
		
		\vspace{0.3cm}
		
		\textbf{Moteur à piston}
		
		\[ T= 	\frac{n_p(550P_{BHP})	}{	V_R/\sqrt{2}	}	\]
		
		\vspace{0.3cm}
		
		\textbf{Moteur à turbopropulseur}
		
		\[	T=T\left(\frac{V_R}{\sqrt{2}}\right)	\]
		
		\vspace{0.3cm}
		
		Un cas particulier de cette méthode consiste à déterminer la distance de décollage pour un train tricycle.
		
		Dans ce cas :
		
		\begin{equation*}
			S_G=	\frac{	V_R^2 W	}{	n_p 550 P_H		+	16.09\rho V_R^3 S 	(\mu C_{LTO}-C_{DTO})	-	2g\mu W V_R		}
		\end{equation*}
		
		%============================================================================================================
		%
		%================================================================================================================
		\subsubsection{ \textcolor{green} { Méthode d’intégration numérique}}
		
		\textbf{\textcolor{yellow}{Avantages}}
		
		Cette méthode permet de modéliser les cas de décollage les plus complexes, notamment les distances de freinage en cas de panne moteur.
		
		\vspace{0.3cm}
		
		\textbf{\textcolor{yellow}{Approximation}}
		
		La vitesse de décollage est approximée par :	\[	V_{LOF}=1.1V_s	\]	où $V_s$ est la vitesse de décrochage.
		
		\vspace{0.3cm}
		

		
		La méthode numérique utilisée pour calculer la distance de roulement repose sur une intégration pas à pas des équations du mouvement.
		
		\paragraph{Étape 1 : Calcul des forces}
		
		À chaque pas de temps $i$, calculer :
		
		\begin{itemize}
			\item la poussée : $T_i$
			\item la traînée : $D_i$
			\item la portance : $L_i$
		\end{itemize}
		
		La pression dynamique :		\[ q_i=\frac{1}{2}\rho V_{i-1}^2 \] où :
		
		\begin{itemize}
			\item $\rho$ est la densité de l’air ;
			\item $V_{i-1}$ est la vitesse au pas précédent.
		\end{itemize}
		
		\paragraph{Étape 2 : Calcul de l’accélération}
		
		\[ a_i= \frac{g}{W}	\left[	T_i-D_i-\mu(W-L_i)	\right]	\] 	avec :		
		\begin{itemize}
			\item $g$ : accélération gravitationnelle ;
			\item $W$ : poids de l’avion ;
			\item $\mu$ : coefficient de frottement.
		\end{itemize}
		
		\paragraph{Étape 3 : Mise à jour de la vitesse et de la distance}
		
		\[	V_i=V_{i-1}+a_i\Delta t_i	\]
		
		\[	S_i=S_{i-1}	+ V_{i-1}\Delta t_i + \frac{1}{2}a_i\Delta t_i^2 \]		avec :	\[	\Delta t_i=t_i-t_{i-1}	\]
		
		\paragraph{Étape 4 : Condition de décollage}
		
		Le calcul est répété jusqu’à :	\[	V_i=V_{LOF}	\]
		
		La distance obtenue :		\[	S_i=S_G	\]
		
		représente la distance totale de roulement au sol.
		
		
		
		
		
		
		
		
		
		
		
		
		%------------------------------------------------
